\documentclass{article}

\def\npart {IA}
\def\nterm {Lent}
\def\nyear {2024}
\def\nlecturer {Prof. C.\ Warnick}
\def\ncourse {Probability}

\setcounter{tocdepth}{2}

\usepackage{alltt}
\usepackage{amsmath, amsfonts, amssymb, amsthm}
\usepackage{esdiff}
\usepackage{booktabs}
\usepackage{caption}
\usepackage{enumitem}
\usepackage{fancyhdr}
\usepackage{graphicx}
\usepackage{mathdots}
\usepackage{mathtools}
\usepackage{microtype}
\usepackage{multirow}
\usepackage{pdflscape}
\usepackage{pgfplots}
\usepackage{siunitx}
\usepackage{textcomp}
\usepackage{slashed}
\usepackage{tabularx}
\usepackage{tikz}
\usetikzlibrary{graphs}
\usepackage{tkz-euclide}
\usepackage[normalem]{ulem}
\usepackage[all]{xy}

\pgfplotsset{compat=1.18}
\pagestyle{fancyplain}

\lhead{\emph{\nouppercase{\leftmark}}}
\rhead{
  \ifnum\thepage=1
  \else
    \npart\ \vline\ \ncourse
  \fi}

\def\nauthor{Marcus Ng}
\author{Based on lectures by \nlecturer \\\small Notes taken by \nauthor}
\date{\nterm\ \nyear}
\title{Part \npart\ --- \ncourse}

\newcommand\numberthis{\addtocounter{equation}{1}\tag{\theequation}}
\newcommand\ilnumberthis{\addtocounter{equation}{1}(\theequation)}

\newcommand*{\Cdot}{{\raisebox{-0.25ex}{\scalebox{1.5}{$\cdot$}}}}
\setlist[enumerate,1]{label={(\roman*)}}

\newcounter{temp_equation_var}
\SetEnumitemKey{eqn}{
  before={
    % Work around usecounter behaviour that sets the counter to 0
    \setcounter{temp_equation_var}{\value{equation}}
    \usecounter{equation}
    \setcounter{equation}{\value{temp_equation_var}}
  },
  label=(\theequation),
  ref={\theequation},
  align=left,
  leftmargin=\parindent
}

% Theorems
\theoremstyle{definition}
\newtheorem*{aim}{Aim}
\newtheorem*{axiom}{Axiom}
\newtheorem*{claim}{Claim}
\newtheorem*{cor}{Corollary}
\newtheorem*{conjecture}{Conjecture}
\newtheorem*{defi}{Definition}
\newtheorem*{eg}{Example}
\newtheorem*{ex}{Exercise}
\newtheorem*{fact}{Fact}
\newtheorem*{law}{Law}
\newtheorem*{lemma}{Lemma}
\newtheorem*{notation}{Notation}
\newtheorem*{prop}{Proposition}
\newtheorem*{question}{Question}
\newtheorem*{problem}{Problem}
\newtheorem*{rrule}{Rule}
\newtheorem*{thm}{Theorem}
\newtheorem*{assumption}{Assumption}
\newtheorem*{assert}{Assertion}

\newtheorem*{remark}{Remark}
\newtheorem*{warning}{Warning}

\newtheorem{nthm}{Theorem}[section]
\newtheorem{nlemma}[nthm]{Lemma}
\newtheorem{nprop}[nthm]{Proposition}
\newtheorem{ncor}[nthm]{Corollary}

\SetEnumitemKey{cases}{
  label=\underline{Case~\arabic*:},
  ref={\arabic*},
  align=left,
  leftmargin=\parindent
}

% \renewcommand{\labelitemi}{--}
% \renewcommand{\labelitemii}{$\circ$}
% \renewcommand{\labelenumi}{(\roman{*})}

\let\stdsection\section
\renewcommand\section{\newpage\stdsection}

% Strike through
\def\st{\bgroup \ULdepth=-.55ex \ULset}

%%%%%%%%%%%%%%%%%%%%%%%%%
%%%%% Maths Symbols %%%%%
%%%%%%%%%%%%%%%%%%%%%%%%%

% Logical symbols
\newcommand{\contradiction}{
  \begin{tikzpicture}[rotate=45,x=0.5ex,y=0.5ex]
  \draw[line width=.1ex] (0,2) -- (3,2) (0,1) -- (3,1) (1,3) -- (1,0) (2,3) -- (2,0);
  \end{tikzpicture}
}

\let\U\relax
\let\C\relax
\let\G\relax

% Matrix groups
% \newcommand{\GL}{\mathrm{GL}}
% \newcommand{\Or}{\mathrm{O}}
% \newcommand{\PGL}{\mathrm{PGL}}
% \newcommand{\PSL}{\mathrm{PSL}}
% \newcommand{\PSO}{\mathrm{PSO}}
% \newcommand{\PSU}{\mathrm{PSU}}
% \newcommand{\SL}{\mathrm{SL}}
% \newcommand{\SO}{\mathrm{SO}}
% \newcommand{\Spin}{\mathrm{Spin}}
% \newcommand{\Sp}{\mathrm{Sp}}
% \newcommand{\SU}{\mathrm{SU}}
\newcommand{\U}{\mathrm{U}}
% \newcommand{\Mat}{\mathrm{Mat}}

% Matrix algebras
% \newcommand{\gl}{\mathfrak{gl}}
% \newcommand{\ort}{\mathfrak{o}}
% \newcommand{\so}{\mathfrak{so}}
% \newcommand{\su}{\mathfrak{su}}
% \newcommand{\uu}{\mathfrak{u}}
% \renewcommand{\sl}{\mathfrak{sl}}

% Calculus
\newcommand*\diffd{\mathop{}\!\mathrm{d}}
\newcommand{\ltlo}[1]{\underline{o}\bkt{#1}}

% Special sets
\newcommand{\C}{\mathbb{C}}
\newcommand{\CP}{\mathbb{CP}}
\newcommand{\GG}{\mathbb{G}}
\newcommand{\N}{\mathbb{N}}
\newcommand{\Q}{\mathbb{Q}}
\newcommand{\R}{\mathbb{R}}
\newcommand{\RP}{\mathbb{RP}}
\newcommand{\T}{\mathbb{T}}
\newcommand{\Z}{\mathbb{Z}}
\renewcommand{\H}{\mathbb{H}}

% Brackets
\newcommand{\bkt}[1]{\left( #1 \right)}
\newcommand{\sqbkt}[1]{\left[ #1 \right]}
\newcommand{\crbkt}[1]{\left\{ #1 \right\}}
\newcommand{\abs}[1]{\left\lvert #1\right\rvert}
\newcommand{\anglebk}[1]{\left\langle #1\right\rangle}
\newcommand{\bket}[1]{\left\lvert #1\right\rangle}
\newcommand{\brak}[1]{\left\langle #1 \right\rvert}
\newcommand{\braket}[2]{\left\langle #1\middle\vert #2 \right\rangle}
\newcommand{\bra}{\langle}
\newcommand{\ket}{\rangle}
\newcommand{\norm}[1]{\left\lVert #1\right\rVert}
\newcommand{\normalorder}[1]{\mathop{:}\nolimits\!#1\!\mathop{:}\nolimits}
\newcommand{\tv}[1]{|#1|}

\newcommand{\pvec}[1]{\overrightarrow{#1}}
\renewcommand{\vec}[1]{\underline{#1}}
\newcommand{\uvec}[1]{\hat{\vec{#1}}}

% not-math
% \newcommand{\bolds}[1]{{\bfseries #1}}
% \newcommand{\cat}[1]{\mathsf{#1}}
% \newcommand{\ph}{\,\cdot\,}
% \newcommand{\term}[1]{\emph{#1}\index{#1}}
% \newcommand{\phantomeq}{\hphantom{{}={}}}

% Probability
% \DeclareMathOperator{\Bernoulli}{Bernoulli}
% \DeclareMathOperator{\betaD}{beta}
% \DeclareMathOperator{\bias}{bias}
% \DeclareMathOperator{\binomial}{binomial}
% \DeclareMathOperator{\corr}{corr}
% \DeclareMathOperator{\cov}{cov}
% \DeclareMathOperator{\gammaD}{gamma}
% \DeclareMathOperator{\mse}{mse}
% \DeclareMathOperator{\multinomial}{multinomial}
% \DeclareMathOperator{\Poisson}{Poisson}
% \DeclareMathOperator{\var}{var}
% \newcommand{\E}{\mathbb{E}}
% \newcommand{\Prob}{\mathbb{P}}

% Algebra
% \DeclareMathOperator{\adj}{adj}
% \DeclareMathOperator{\Ann}{Ann}
% \DeclareMathOperator{\Aut}{Aut}
% \DeclareMathOperator{\Char}{char}
% \DeclareMathOperator{\disc}{disc}
% \DeclareMathOperator{\dom}{dom}
% \DeclareMathOperator{\fix}{fix}
\DeclareMathOperator{\Hom}{Hom}
% \DeclareMathOperator{\id}{id}
% \DeclareMathOperator{\image}{im}
\DeclareMathOperator{\im}{im}
% \DeclareMathOperator{\re}{re}
% \DeclareMathOperator{\tr}{tr}
% \DeclareMathOperator{\Tr}{Tr}
% \newcommand{\Bilin}{\mathrm{Bilin}}
% \newcommand{\Frob}{\mathrm{Frob}}

% Others
% \newcommand\ad{\mathrm{ad}}
% \newcommand\Art{\mathrm{Art}}
% \newcommand{\B}{\mathcal{B}}
% \newcommand{\cU}{\mathcal{U}}
% \newcommand{\Der}{\mathrm{Der}}
\newcommand{\D}{\mathrm{D}}
% \newcommand{\dR}{\mathrm{dR}}
% \newcommand{\exterior}{\mathchoice{{\textstyle\bigwedge}}{{\bigwedge}}{{\textstyle\wedge}}{{\scriptstyle\wedge}}}
% \newcommand{\F}{\mathbb{F}}
% \newcommand{\G}{\mathcal{G}}
% \newcommand{\Gr}{\mathrm{Gr}}
% \newcommand{\haut}{\mathrm{ht}}
% \newcommand{\Hol}{\mathrm{Hol}}
% \newcommand{\hol}{\mathfrak{hol}}
\newcommand{\Id}{\mathrm{Id}}
% \newcommand{\lie}[1]{\mathfrak{#1}}
% \newcommand{\op}{\mathrm{op}}
% \newcommand{\Oc}{\mathcal{O}}
% \newcommand{\pr}{\mathrm{pr}}
% \newcommand{\Ps}{\mathcal{P}}
% \newcommand{\pt}{\mathrm{pt}}
% \newcommand{\qeq}{\mathrel{``{=}"}}
% \newcommand{\Rs}{\mathcal{R}}
% \newcommand{\Vect}{\mathrm{Vect}}
% \newcommand{\wsto}{\stackrel{\mathrm{w}^*}{\to}}
% \newcommand{\wt}{\mathrm{wt}}
% \newcommand{\wto}{\stackrel{\mathrm{w}}{\to}}
% \renewcommand{\d}{\mathrm{d}}
% \renewcommand{\P}{\mathbb{P}}
%\renewcommand{\F}{\mathcal{F}}

\let\Im\relax
\let\Re\relax

% \DeclareMathOperator{\area}{area}
% \DeclareMathOperator{\card}{card}
% \DeclareMathOperator{\ccl}{ccl}
% \DeclareMathOperator{\ch}{ch}
% \DeclareMathOperator{\cl}{cl}
% \DeclareMathOperator{\cls}{\overline{\mathrm{span}}}
% \DeclareMathOperator{\coker}{coker}
% \DeclareMathOperator{\conv}{conv}
% \DeclareMathOperator{\cosec}{cosec}
% \DeclareMathOperator{\cosech}{cosech}
% \DeclareMathOperator{\covol}{covol}
% \DeclareMathOperator{\diag}{diag}
% \DeclareMathOperator{\diam}{diam}
% \DeclareMathOperator{\End}{End}
% \DeclareMathOperator{\energy}{energy}
% \DeclareMathOperator{\erfc}{erfc}
% \DeclareMathOperator{\erf}{erf}
% \DeclareMathOperator*{\esssup}{ess\,sup}
% \DeclareMathOperator{\ev}{ev}
% \DeclareMathOperator{\Ext}{Ext}
% \DeclareMathOperator{\fst}{fst}
% \DeclareMathOperator{\Fit}{Fit}
% \DeclareMathOperator{\Frac}{Frac}
% \DeclareMathOperator{\Gal}{Gal}
% \DeclareMathOperator{\gr}{gr}
% \DeclareMathOperator{\hcf}{hcf}
\DeclareMathOperator{\Im}{Im}
% \DeclareMathOperator{\Ind}{Ind}
% \DeclareMathOperator{\Int}{Int}
\DeclareMathOperator{\Isom}{Isom}
% \DeclareMathOperator{\lcm}{lcm}
% \DeclareMathOperator{\length}{length}
% \DeclareMathOperator{\Lie}{Lie}
% \DeclareMathOperator{\like}{like}
% \DeclareMathOperator{\Lk}{Lk}
% \DeclareMathOperator{\Maps}{Maps}
% \DeclareMathOperator{\orb}{orb}
% \DeclareMathOperator{\ord}{ord}
% \DeclareMathOperator{\otp}{otp}
% \DeclareMathOperator{\poly}{poly}
% \DeclareMathOperator{\rank}{rank}
% \DeclareMathOperator{\rel}{rel}
% \DeclareMathOperator{\Rad}{Rad}
\DeclareMathOperator{\Re}{Re}
% \DeclareMathOperator*{\res}{res}
% \DeclareMathOperator{\Res}{Res}
% \DeclareMathOperator{\Ric}{Ric}
% \DeclareMathOperator{\rk}{rk}
% \DeclareMathOperator{\Rees}{Rees}
% \DeclareMathOperator{\Root}{Root}
% \DeclareMathOperator{\sech}{sech}
% \DeclareMathOperator{\sgn}{sgn}
% \DeclareMathOperator{\snd}{snd}
% \DeclareMathOperator{\Spec}{Spec}
\DeclareMathOperator{\spn}{span}
% \DeclareMathOperator{\stab}{stab}
% \DeclareMathOperator{\St}{St}
% \DeclareMathOperator{\supp}{supp}
% \DeclareMathOperator{\Syl}{Syl}
\DeclareMathOperator{\Sym}{Sym}
% \DeclareMathOperator{\vol}{vol}
\numberwithin{equation}{section}

\begin{document}
\maketitle{
    \small
    \noindent\textbf{Limits and Convergence}\\
    Sequences and series in $\R$ and $\C$. Sums, products and quotients.
    Absolute convergence; absolute convergence implies convergence. 
    The Bolzano-Weierstrass theorem and applications (the General Principle of Convergence). 
    Comparison and ratio tests, alternating series test.\hspace*{\fill}[6]
    
    \vspace{10pt}
    \noindent\textbf{Continuity}\\
    Continuity of real and complex valued functions defined on subsets of $\R$ and $\C$.
    The intermediate value theorem. A continuous function on a closed bounded interval is bounded and attains its bounds.\hspace*{\fill}[5]

    \vspace{10pt}
    \noindent\textbf{Differentiability}\\
    Differentiability of functions from $\R$ to $\R$. Derivative of sums and products. The chain rule.
    Derivative of the inverse function. Roll's theorem; the mean value theorem. One-dimensional version of the inverse function theorem.
    Taylor's theorem from $\R$ to $\R$; Lagrange's form of the remainder. Complex differentiaion.\hspace*{\fill}[5]
        
    \vspace{10pt}
    \noindent\textbf{Power series}\\
    Complex power series and radius of convergence. Exponential, trigonometric and hyperbolic functions, and relations between them.
    *Direct proof of the differentiability of a power series within its circle of convergence*.\hspace*{\fill}[4]

    \vspace{10pt}
    \noindent\textbf{Integration}\\
    Definition and basic properties of the Riemann integral. A non-integrable function.
    Integrability of monotonic functions. Integrability of piecewise-continuous functions.
    The fundamental theorem of calculus. Differentiation of indefinite integrals.
    Integration by parts. The integral form of the remainder in Taylor's theorem. Impropper integrals.\hspace*{\fill}[6]
}
\tableofcontents

%%%%%%%%%%%%%%%%%%%%%%
% Lecture 1: 19/01/2024
%%%%%%%%%%%%%%%%%%%%%%

\section{Limits and Convergence}
Consider a sequence of real numbers $a_1, a_2, a_3, \cdots$ which we write as $(a_n)$ or $(a_n)_{n=1}^\infty$ for $a_n \in \R$.

\begin{defi}[Limit]
    We say that $a_n \rightarrow a$ as $n \rightarrow \infty$ for some $a \in \R$ if:
    Given $\epsilon > 0$, $\exists N \in \N$ such that $\forall n \geq N$
    \[\abs{a_n - a} < \epsilon\]
\end{defi}

\begin{remark}
    We can freely replace the strict inequality with $\leq$ and it does not change the definition. 
    We can also replace $\epsilon$ with any multiple of $\epsilon$ say $5 \epsilon$.
    We also assume that $\infty \notin \R$, so $a = \infty$ is not allowed.
\end{remark}
\begin{defi}[Increasing and decreasing]\leavevmode
    \begin{enumerate}
        \item If $a_n \leq a_{n+1}, \ \forall n$, then $(a_n)$ is \emph{increasing}
        \item If $a_n \geq a_{n+1}, \ \forall n$, then $(a_n)$ is \emph{decreasing} 
        \item If $a_n < a_{n+1}, \ \forall n$, then $(a_n)$ is \emph{strictly increasing} 
        \item If $a_n > a_{n+1}, \ \forall n$, then $(a_n)$ is \emph{strictly decreasing} 
    \end{enumerate}
\end{defi}

In all cases $(a_n)$ is \emph{monotone}. This allows us to state the \emph{fundamental axiom of real numbers}.

\begin{thm}[Fundamental axiom of real numbers]
    An increasing sequence of real numbers which is bounded above converges (i.e, it has a limit).
    For instance, if $a_n \in \R$ with
    \[a_1 \leq a_2 \leq \cdots \ \text{ and } \ a_n \leq A \ \forall n\]
    There exists $a \in \R$ such that $a_n \rightarrow a$ as $n \rightarrow \infty$.
    Equivalently, a decreasing sequence bounded below has a limit.
\end{thm}

This is equivalent to the \emph{least upper bound axiom}, that states that any non-empty set that is bounded above has a \emph{supremum}.
\begin{defi}[Supremum]
    If $S \subset R$ and $S \neq \emptyset$ we say that
    \[\sup(S) = k\]
    if 
    \begin{enumerate}
        \item $x \leq k \ \forall x \in S$ ($k$ is an upperbound for $S$)
        \item Given $\epsilon > 0 \ \exists x \in S$ such that $x > k - \epsilon$ ($k$ is a least upper bound)
    \end{enumerate}
\end{defi}

A similar definition for greatest lower bound (infimum). 
If $\sup(S) \in S$ we say that $\sup(S)$ is the minimum of $S$ and write
\[\sup(S) = \max(S)\]

\begin{ex}
    Show that the fundamental axiom is equivalent to the least upperbound axiom.
\end{ex}

\begin{nlemma}\leavevmode
    \begin{enumerate}
        \item The limit is unique, i.e; if $a_n \rightarrow a$ and $a_n \rightarrow b$ as $n \rightarrow \infty$ then $a = b$.
        \item If $a_n \rightarrow a$ as $n \rightarrow \infty$ and $n_1 < n_2 < n_3 < \cdots$ is a set of increasing indicies, 
        then $a_{n_j} \rightarrow a$ as $j \rightarrow \infty$. (subsequences converge to the same limit)
        \item If $a_n = c \ \forall n$ then $a_n \rightarrow c$ as $n \rightarrow \infty$
        \item If $a_n \rightarrow a$ and $b_n \rightarrow b$ then $a_n + b_n \rightarrow a + b$ 
        \item If $a_n \rightarrow a$ and $b_n \rightarrow b$ then $a_nb_n \rightarrow ab$
        \item If $a_n \rightarrow a$ and $a_n \neq 0$ and $a \neq 0$ then $\frac{1}{a_n} \rightarrow \frac{1}{a}$
        \item If $a_n \leq A$ and $a_n \rightarrow a$ then $a \leq A$. (note that this is not true if we insist $a < A$)
        \item If $a_n \rightarrow a$ and $c_n \rightarrow a$ as $n \rightarrow \infty$ and
        \[a_n \leq b_n \leq c_n \rightarrow b_n \rightarrow a\]
    \end{enumerate}
\end{nlemma}

\begin{proof}
    We will do (i), (ii) and (v), the others wil be exercises.
    \begin{enumerate}
        \item For any $\epsilon >0$ find $N_1(\epsilon)$ and $N_2(\epsilon)$ such that
        \begin{align*}
            n \geq N_1(\epsilon) \Rightarrow \abs{a_n - a} < \epsilon
            n \geq N_2(\epsilon) \Rightarrow \abs{a_n - b} < \epsilon
        \end{align*}
        Hence if $n \geq \max(N_1, N_2)$ then,
        \begin{align*}
            0 \leq \abs{b - a} &= \abs{b - a_n + a_n - a} \\
            &\leq \abs{b - a_n} + \abs{a_n - a} \\
            &< 2\epsilon
        \end{align*}

        This implies that $\abs{b - a} = 0$, as otherwise we can set $\epsilon = \frac{1}{3} \abs{b - a}$ to find
        \[0 \leq \abs{b - a} < \frac{2}{3} \abs{b - a} \ \contradiction\]

        \item Since $n_j < n_{j + 1} \Rightarrow n_{j+1} \geq n_j + 1$, and hence by induction we have $n_j \geq j$.
        As $a_n \rightarrow a$ as $n \rightarrow \infty$, given $\epsilon >0 \ \exists N = N(\epsilon)$ such that
        \[n \geq N(\epsilon) \Rightarrow \abs{a_n - a} < \epsilon\]
        So if we choose $j \geq N(\epsilon)$ then $n_j \geq j \geq N(\epsilon)$ which implies that
        \[\abs{a_{n_j} - a} < \epsilon\]
        Hence $a_{n_j} \rightarrow a$ as $j \rightarrow \infty$
        
        \item As $a_n \rightarrow a$ and $b_n \rightarrow b$ as $n \rightarrow \infty$ for any $\epsilon > 0 \ \exists N_1 = N_1(\epsilon), N_2 = N_2(\epsilon)$
        such that
        \begin{align*}
            n \geq N_1(\epsilon) \Rightarrow \abs{a_n - a} < \epsilon \\
            n \geq N_2(\epsilon) \Rightarrow \abs{b_n - a} < \epsilon
        \end{align*}
        Now,
        \begin{align*}
            \abs{a_nb_n - ab} &= \abs{a_nb_n - a_nb + a_nb - ab} \\
            &\leq \abs{a_nb_n - a_nb} + \abs{a_nb - ab} \\
            &= \abs{a_n}\abs{b_n - b} + \abs{b}{a_n -a}
        \end{align*}
        If $n \geq N_1$ then $\abs{a_n - a} \leq 1$, therefore
        \[\abs{a_n} = \abs{a_n + a - a} \leq \abs{a_n - a} + \abs{a} \leq (1 + \abs{a})\]
        Thus if $n \geq N_3 = \max(N_1, N_2)$ then,
        \[\abs{a_nb_n - ab} < (1 + \abs{a})\epsilon + \abs{b}\epsilon = (1 + \abs{a} + \abs{b})\epsilon\]
        which can be made arbitrarily small. 
        Therefore
        \[a_nb_n \rightarrow ab\]
    \end{enumerate}
\end{proof}

\end{document}